\documentclass[11pt,reqno]{amsart}

% packages
\usepackage{amsmath, amsthm, amssymb}
\usepackage{setspace}
\setstretch{1.15}
\usepackage{mathtools}
\usepackage{geometry}
\geometry{margin=1in}
\usepackage{hyperref}
\usepackage{enumitem}
\usepackage{graphicx}

\title{A Divisibility Condition for Moduli in Disjoint Covering Systems with Precisely One Repeated Modulus}
\author{Yu Hashimoto}
\date{June 2026}

% theorem environments
\newtheorem{theorem}{Theorem}[section]
\newtheorem{lemma}[theorem]{Lemma}
\newtheorem{corollary}[theorem]{Corollary}
\newtheorem{proposition}[theorem]{Proposition}
\theoremstyle{definition}
\newtheorem{definition}[theorem]{Definition}
\newtheorem{remark}[theorem]{Remark}
\newtheorem{problem}[theorem]{Problem}
\newtheorem{question}[theorem]{Question}
% \theoremstyle{example}
% \newtheorem{example}[theorem]{Example}

% macros
\newcommand{\Zmod}[1]{\mathbb{Z}/#1\mathbb{Z}}

% ind as operator
\DeclareMathOperator{\ind}{index}

% families of cosets
\newcommand{\Cset}{\mathcal{C}}
\newcommand{\Aset}{\mathcal{A}}
\newcommand{\Lset}{\mathcal{L}}
\newcommand{\Dset}{\mathcal{D}}

% show only referenced equation numbers
\mathtoolsset{showonlyrefs=true}

\begin{document}

\maketitle

\section{Introduction}
In this note, we give a proof of the divisibility condition for moduli in a \emph{disjoint covering system} (DCS) with precisely one repeated modulus. This condition was assumed in the author's earlier paper \cite{Hashimoto}.
In \cite{Hashimoto}, we assumed the general theorem on moduli in DCS which was proved by Porubsk\'y \cite{Porubsky} and the Berger, Felzenbaum, and Fraenkel's argument in \cite{BFF1988}. We give the alternative proof for completeness and Lean formalization. The essentially same approach was found in the literature in the form of a generating function \cite{NovakZnam} \cite{Porubsky}.

\section{The proof}

\begin{proposition}
\label{prop:weak-completion-divisibility}
Let \((a_i+n_i\mathbb{Z})_{i=1}^t\)
be a DCS satisfying 
\begin{equation}\label{eq:distinctas}
n_1 < n_2 < \cdots < n_{t-m} < n_{t-m+1}=\cdots=n_t
\qquad (m\ge 2).
\end{equation}
Then
\[
n_i \mid n_t
\qquad
(1\le i\le t).
\]
In particular,
\[
N:=\operatorname{lcm}(n_1,\dots,n_t) \mid n_t.
\]
\end{proposition}

\begin{proof}
Let $1 \le r \le t$ and let \(\zeta\) be a primitive \(n_r\)-th root of unity.
For a DCS, we have the standard identity
\[
\sum_{x\in\Zmod N}\zeta^x
=
\sum_{k=1}^t
\sum_{\substack{x\in\Zmod N\\x\equiv a_k\pmod{n_k}}}
\zeta^x.
\]
since each $x$ is covered exactly once.
For each progression,
\[
\sum_{\substack{
  x \in \Zmod{N} \\
  x \equiv a_k \pmod{n_k}
}}
\zeta^x
=
\zeta^{a_k} \sum_{j = 0}^{N/n_k - 1}(\zeta^{n_k})^{j}
= 
\begin{cases}
  \frac{N}{n_k}\zeta^{a_k} & (n_r \mid n_k) \\
  0 & (n_r \nmid n_k)
\end{cases}
.
\]
Hence 
\[
\sum_{k \colon n_r \mid n_k} \frac{N}{n_k}\zeta^{a_k}
=
0.
\]
Now suppose $n_i \nmid n_t$ for some $1 \le i < t$. Choose $n_r$ maximal, with respect to divisibility, among
the moduli $n_j$ satisfying $n_j\nmid n_t$.
It follows that $n_r \nmid n_k$ for all $k \neq r$. Indeed, if $n_r \mid n_k$ for some $k \neq r$, then $n_k \nmid n_t$ by the choice of $n_r$. Since $k \neq r$ and moduli smaller than $n_t$ are pairwise distinct, $n_r \mid n_k$ contradicts the maximality of $n_r$. Therefore, the above sum has only the $k = r$ term and 
\[
\frac{N}{n_r}\zeta^{a_r} = 0.
\]
This is a contradiction and hence completes the proof.
\end{proof}

\bibliographystyle{plain}
\bibliography{references}

\end{document}
